\begin{scanpage}{0174}{156}
the screen, the author obtained an approximation solution in the form of a surface integral. The boundary values of the integrand were taken from the rigorous solution to the problem of diffraction by a half-plane which was found by Sommerfeld. The field was calculated in the far zone on the axis of the hole and far from it, and also on the axis near the screen. Using this approach, Braunbek recently solved the problem of scalar wave diffraction by an aperture in a concially shaped screen [31].

In the papers of Frahn [32, 33], this method was used for the diffraction of electromagnetic waves. Diffraction of a plane wave incident normal to an ideally conducting screen with a circular hole was investigated. The field was calculated in the hole and on the axis, and also the field in the far zone and the transmission coefficient (the ratio of the energy passing through the hole to the energy falling on it) were calculated.

In these works of Braunbek and Frahn, secondary diffraction was not considered. The expressions obtained by them for the fringing field intensity in the far zone agree with similar expressions following from our equations (\S~9).

Karp and Russek [51] studied diffraction by a slit in the case when the incident wave's electric vector is parallel to the slit edge. They investigated each semi-infinite part of the screen as a half-plane excited by the incident wave field and a ``virtual'' source localized on the edge of the opposite half-plane. The moments of these sources were determined from a system of two algebraic equations which were obtained by using the asymptotic expressions resulting from the rigorous solution for the half-plane. Secondary diffraction was considered, and partially the general interaction. Special attention was allotted to calculating the transmission coefficient, but equations for the scattering characteristics which would be suitable with all directions of incident wave propagation were absent.
\end{scanpage}
